% \sham: plain, full, and forgotten after a source whose heading has no reference.
\documentclass{article}
\usepackage{shiur}
\setshiurfont{hebrew}{FreeSerif}
\begin{document}
\begin{mekor}
\begin{mekorlang}{english}{\bavli{menachot}{69b}}First.\end{mekorlang}
\end{mekor}
\begin{mekor}
\begin{mekorlang}{english}{Rashi \sham}Second.\end{mekorlang}
\end{mekor}
\begin{mekor}
\begin{mekorlang}{english}{Bereishit Rabbah 87:7}Third.\end{mekorlang}
\end{mekor}
\mareimekomotsetup{sham=full}
\begin{mekor}
\begin{mekorlang}{english}{Tosafot \sham}Fourth.\end{mekorlang}
\end{mekor}
\begin{mekor}
\begin{mekorlang}{english}{\tanakh[rashi]{gen}{18:7}}Fifth.\end{mekorlang}
\end{mekor}
\begin{mekor}
\begin{mekorlang}{english}{Ibn Ezra \sham}Sixth.\end{mekorlang}
\end{mekor}
\begin{mekor}
\begin{mekorlang}{hebrew}{רמב״ן \sham}שביעי.\end{mekorlang}
\end{mekor}
\end{document}

%! expect: Rashi ad loc Second.
%! expect: Tosafot ad loc Fourth.
%! expect: Ibn Ezra ad loc (Genesis 18:7) Sixth.
% (pdftotext scrambles the order of a bracketed reference in Hebrew, so its parts
% are checked separately)
%! expect: רמב״ן שם
%! expect: בראשית י״ח
